\section{Related Work}
\label{sec:relatedwork}

Our work lies at the intersection of (i)~learning-based respiratory sound analysis, (ii)~edge/IoT deployment constraints, (iii)~dynamic inference and model compression, and (iv)~heterogeneous FPGA acceleration. This section reviews relevant literature along these axes and clarifies the gap addressed by the proposed UA-HCI framework.

\subsection{Learning-Based Respiratory Sound Classification}
The ICBHI 2017 challenge and its publicly released respiratory sound database~\cite{b16} catalyzed the transition from handcrafted acoustic descriptors toward deep representation learning for automated lung sound analysis. Since the challenge, a range of deep architectures have been explored. Convolutional--recurrent designs combine a spectral CNN encoder with sequence modeling to capture both time--frequency patterns and temporal dynamics in adventitious sounds~\cite{noureddine2020}. More recent work has investigated attention-centric approaches; for example, Ariyanti \textit{et al.} proposed an audio-spectrogram vision transformer for abnormal respiratory sound identification on ICBHI~2017~\cite{b07}. Other studies have explored CNN-based disease classification using MFCCs and spectrogram variants~\cite{b08,b10,b12,b43,b45} as well as lightweight architectures designed to reduce model complexity for embedded deployment~\cite{b13,b14}. Multi-task learning has also been explored for simultaneous sound and disease classification~\cite{b15}.

A recurring issue in this literature is the sensitivity of reported accuracy to evaluation protocol. Many studies use random segment-level splits, whereas subject-/patient-independent splitting is more conservative for estimating clinical generalization~\cite{b16}. When comparing methods, it is therefore important to report the splitting strategy, label mapping, and unit of inference.

\subsection{Dynamic Inference and Early-Exit Strategies}
Static inference applies the same computational cost to every sample regardless of difficulty. \emph{Dynamic inference} methods adapt computation to input complexity by enabling early termination when model confidence is high. BranchyNet~\cite{branchynet} is a representative approach that attaches auxiliary classifiers at intermediate network depths, allowing fast inference via early exits for easy inputs. Multi-scale dense networks~\cite{msdnet2018} extend this idea by providing multiple classifier heads at different feature scales with dense cross-layer connections, trading computation for accuracy depending on the exit point. The Once-for-All framework~\cite{ofa2020} trains a single large network from which specialized sub-networks of varying depth and width can be extracted without retraining, supporting resource-adaptive deployment. Although these methods were originally proposed in vision, the underlying principle---allocating compute proportionally to sample difficulty---is directly applicable to edge respiratory analysis where energy and latency are constrained.

\subsection{Model Compression, Quantization, and Knowledge Distillation}
Complementing architectural efficiency, model compression techniques reduce the cost of each inference pass. \emph{Knowledge distillation} transfers predictive behavior from a large teacher model to a smaller student, typically by matching softened output distributions~\cite{hinton2015distilling}. This technique is especially relevant when the deployment target is constrained, as in the present work where an EfficientNet-B0 teacher ensemble~\cite{efficientnet2019} distills into a MobileNetV2 student~\cite{b35}.

\emph{Quantization} reduces the numerical precision of weights and activations from floating-point to low-bit integer representations, enabling efficient execution on hardware accelerators. Foundational work on quantized neural networks demonstrated that low-precision training is feasible with minimal accuracy loss~\cite{hubara2016qnn,jacob2017quantization}. Comprehensive guides to quantization-aware training (QAT) for CNNs~\cite{krishnamoorthi2018quant} and practical 8-bit parallelism strategies~\cite{dettmers2015} have further streamlined deployment on integer-only accelerators such as the Xilinx DPU. In our framework, QAT converts the MobileNetV2 student to INT8 for DPU-compatible inference.

Platform-aware neural architecture search (NAS) methods have also contributed to efficient edge models. MnasNet~\cite{mnasnet2019} incorporates latency into the search objective to find architectures that balance accuracy and speed on mobile devices. Together, these techniques form the model-efficiency stack that enables the Layer-4 CNN stage in UA-HCI.

\subsection{FPGA Acceleration and Heterogeneous Co-Design}
FPGAs offer deterministic latency and favorable energy-per-inference compared to mobile GPUs, making them attractive for embedded medical AI. Recent toolchains (e.g., Xilinx Vitis AI) and accelerator IP (e.g., the DPU) have streamlined CNN deployment on FPGAs. Survey and application studies report increasing adoption of FPGA-based acceleration in medical imaging~\cite{zhu2023} and embedded audio classification~\cite{fazilat2023}. Nevertheless, most FPGA deployments still execute the full deep model for every input. This limits the achievable system-level energy reduction when a large fraction of samples could be classified by lightweight screening stages.

Heterogeneous SoC platforms (e.g., Xilinx Zynq UltraScale+) partition computation between an embedded ARM Processing System (PS) and Programmable Logic (PL). This partitioning is natural for cascaded inference: control-oriented lightweight stages run on the PS, while compute-intensive CNN inference is offloaded to the PL only when needed.

\subsection{Summary of the Research Gap}
Prior respiratory sound work demonstrates strong deep-learning performance, and the hardware literature demonstrates efficient FPGA-based CNN acceleration. However, fewer studies jointly address \emph{(i)}~uncertainty-aware routing across lightweight and deep representations, \emph{(ii)}~heterogeneous PS--PL co-design that reserves programmable-logic execution for uncertain samples, \emph{(iii)}~a model-efficiency stack combining distillation, QAT, and efficient architectures, and \emph{(iv)}~route-level measurement protocols that quantify end-to-end latency and energy under dynamic inference. The UA-HCI framework targets this gap by integrating a Random Forest screening cascade on the PS with selective DPU-based MobileNetV2 inference on the PL, supported by QAT and knowledge distillation for the CNN stage.